\chapter{Polynomial Fermat and three cubes}\label{ch:polyfermat}

Polynomials have a derivative, and the derivative proves a Fermat theorem. The integers do not, and the same shape of statement is a conjecture. Both belong with the arithmetic of this volume: one is a theorem about $\mathbb C[t]$, the other is a congruence, two polynomial identities, and a local density, with Heath-Brown's prediction labelled as a prediction.

\section{Polynomials}

For polynomials over $\mathbb C$, the radical becomes the number of distinct roots. For a nonzero polynomial $f$, let $n_0(f)$ be the number of distinct
complex roots of $f$.

\begin{theorem}[Stothers, 1981; Mason, 1984]\label{conjectures:thm:mason}
Let $a,b,c\in\mathbb C[t]$ be coprime polynomials with $a+b=c$, not all constant. Then
\[
\max(\deg a,\deg b,\deg c)\le n_0(abc)-1.
\]
\end{theorem}
\begin{proof}
(Snyder's proof.) Since $a$ and $b$ are coprime and $c=a+b$, the three polynomials are pairwise coprime, so they have no common roots. Consider the
\emph{Wronskian} $W=ab'-a'b$. Substituting $b=c-a$ and $a=c-b$,
\[
W=ac'-a'c=-(bc'-b'c).
\]
\emph{$W\ne0$.} Otherwise $(b/a)'=W/a^2=0$, so $b=\lambda a$ for a constant $\lambda$; as $a$ and $b$ are coprime, both are constant, and so is $c$.

\emph{Many roots of $W$.} Let $\alpha$ be a root of $a$ of multiplicity $m$. Then $(t-\alpha)^{m-1}$ divides $a$ and $a'$, hence $W=ab'-a'b$. The same holds for
the roots of $b$ (same formula) and of $c$ (formula $W=ac'-a'c$). These roots are distinct, so
\[
D=\prod_{\alpha:\ (abc)(\alpha)=0}(t-\alpha)^{m_\alpha-1}
\]
divides $W$, where $m_\alpha$ is the multiplicity of $\alpha$ as a root of $abc$. Its degree is $\deg a+\deg b+\deg c-n_0(abc)$.

\emph{Few roots of $W$.} $\deg W\le\deg a+\deg b-1$: the terms of degree $\deg a+\deg b$ in $ab'$ and $a'b$ vanish (their leading coefficients are
$\deg b\cdot\ell$ and $\deg a\cdot\ell$ times $t^{\deg a+\deg b-1}$ for the product $\ell$ of the leading coefficients, so even degree $\deg a+\deg b-1$ is
an upper bound). Hence
\[
\deg a+\deg b+\deg c-n_0(abc)\le\deg W\le\deg a+\deg b-1,
\]
that is $\deg c\le n_0(abc)-1$. Using $W=ac'-a'c$ (so $\deg W\le\deg a+\deg c-1$) gives $\deg b\le n_0(abc)-1$, and $W=-(bc'-b'c)$ gives
$\deg a\le n_0(abc)-1$.
\end{proof}

\begin{corollary}[Fermat's last theorem for polynomials]\label{conjectures:cor:fermatpoly}
For $n\ge3$ there are no coprime polynomials $x,y,z\in\mathbb C[t]$, not all constant, with $x^n+y^n=z^n$.
\end{corollary}
\begin{proof}
Apply Theorem~\ref{conjectures:thm:mason} to $a=x^n$, $b=y^n$, $c=z^n$, which are coprime. The roots of $x^ny^nz^n$ are the roots of $x$, $y$ and $z$, so
$n_0(x^ny^nz^n)\le\deg x+\deg y+\deg z=:d$. Then $n\deg x\le d-1$, and likewise for $y$ and $z$; adding, $nd\le3d-3<3d$, so $n<3$.
\end{proof}

The bound is sharp for $a=t$, $b=1-t$, $c=1$, up to the coprimality one gets by cancelling a common factor, and the argument uses characteristic $0$ to know that $W\ne0$. It is the model for the
integer conjecture: $\deg$ plays the role of $\log c$, and $n_0$ the role of $\log\operatorname{rad}(abc)$. For polynomials the derivative does the work; the
integers have no derivative, and that is one way to say why abc is hard.


The integer $abc$ conjecture is the same inequality with $\log c$ in place of the degree and $\log\operatorname{rad}(abc)$ in place of the number of roots. It is a conjecture. It is not used anywhere in this volume.

\section{Sums of three cubes}

\section{What can be proved}

\begin{proposition}\label{prop:mod9}
If $n\equiv\pm4\pmod9$, then $n$ is not a sum of three cubes of integers.
\end{proposition}

\begin{proof}
If $3\mid x$ then $x^3\equiv0\pmod{27}$. Otherwise $x\equiv\pm1,\pm2,\pm4\pmod9$, and $x^3\equiv\pm1,\pm8,\pm64$, all $\equiv\pm1\pmod9$. So cubes are $0,\pm1$ modulo $9$, and a sum of three of them is in
$\{0,\pm1,\pm2,\pm3\}$ modulo $9$, which misses $\pm4$.
\end{proof}

For two numbers, infinitely many representations are known, given by polynomial identities.

\begin{proposition}[Verebrusov, 1908; Mahler, 1936]\label{prop:families}
For every integer $t$,
\[
(1+6t^3)^3+(1-6t^3)^3+(-6t^2)^3=2,\qquad(9t^4)^3+(3t-9t^4)^3+(1-9t^3)^3=1.
\]
\end{proposition}

\begin{proof}
For the first, $(1+u)^3+(1-u)^3=2+6u^2$; with $u=6t^3$ this is $2+216t^6=2+(6t^2)^3$. For the second, let $a=9t^4$ and $b=3t-9t^4$, so $a+b=3t$ and
$a^3+b^3=(a+b)^3-3ab(a+b)=27t^3-9t\cdot9t^4(3t-9t^4)=27t^3-243t^6+729t^9$; and $(1-9t^3)^3=1-27t^3+243t^6-729t^9$. The sum is $1$.
\end{proof}

No such family is known for any other $n$ that is not a cube (a cube $n=m^3$ has the family $m^3+t^3+(-t)^3$), and for $3$ it is not even known whether there are infinitely many
solutions. So the conjecture that every $n\not\equiv\pm4\pmod9$ is a sum of three cubes is open; it is believed, because of the heuristic.

\section{The heuristic}

Count solutions in a box. Let $N_n(B)$ be the number of $(x,y,z)\in\Z^3$ with $x^3+y^3+z^3=n$ and $\max(|x|,|y|,|z|)\le B$. The heuristic has two factors, one from the real numbers and one from the
primes; we derive each, and say exactly which step is not a proof.

\subsection*{The real density}

Pretend that the value $F=x^3+y^3+z^3$ at a random integer point near a point of $\R^3$ is ``spread out'' like a continuous variable. The expected number of integer points with $F=n$ is then the
volume of the region where $|F-n|<\frac12$. For a dyadic shell $B<\max(|x|,|y|,|z|)\le2B$, substitute $x=Bu$, $y=Bv$, $z=Bw$: the region becomes $\{1<\max(|u|,|v|,|w|)\le2,\ |u^3+v^3+w^3-n/B^3|<\frac1{2B^3}\}$, of
volume about $\frac1{B^3}A$, where $A$ is the ``area'' of the surface $u^3+v^3+w^3=0$ in the shell, measured with the density $\frac1{|\nabla F|}$; multiplied by $B^3$ (the Jacobian), the shell
contributes a constant $A$, \emph{independent of $B$}. There are $\log_2B$ dyadic shells up to $B$, so the real part of the heuristic predicts
\[
N_n(B)\approx\sigma_\infty\log B,
\]
with an explicit constant $\sigma_\infty$; following Heath-Brown, Booker and Sutherland find $\sigma_\infty=4\int_1^\infty(t^3+1)^{-2/3}\dd t=\frac{2\,\Gamma(\frac13)^2}{3\,\Gamma(\frac23)}\approx3.533$. The logarithm is the signature of the critical degree: three variables of degree three, so each dyadic scale contributes the same amount.

\subsection*{The local densities}

The integers are not random: congruences can forbid solutions, as Proposition~\ref{prop:mod9} shows, or favour them. For a prime $p$, the \emph{local density} measures how much more (or less) often
$F\equiv n$ happens modulo powers of $p$ than for a random value:
\[
\sigma_p(n)=\lim_{k\to\infty}\frac{\#\{(x,y,z)\bmod p^k:\ x^3+y^3+z^3\equiv n\pmod{p^k}\}}{p^{2k}}.
\]
(For a random value, the count would be $p^{3k}\cdot p^{-k}=p^{2k}$.)

\begin{proposition}\label{prop:localdens}
\begin{enumerate}[label=(\alph*)]
\item If $p\equiv2\pmod3$ and $p\nmid n$, then $\sigma_p(n)=1$.
\item If $n\equiv\pm4\pmod9$, then $\sigma_3(n)=0$.
\end{enumerate}
\end{proposition}

\begin{proof}
(a) \emph{Modulo $p$.} The group $\mathbb F_p^\times$ has order $p-1$, prime to $3$, so $x\mapsto x^3$ is injective on it, and since $0^3=0$ it is a bijection of $\mathbb F_p$. Hence $(x,y,z)\mapsto(x^3,y^3,z^3)$ permutes $\mathbb F_p^3$, and
the number of solutions of $x^3+y^3+z^3=n$ in $\mathbb F_p$ equals the number of $(a,b,c)$ with $a+b+c=n$, which is $p^2$.

\emph{Lifting.} Let $F=x^3+y^3+z^3-n$. A solution modulo $p$ has some coordinate nonzero modulo $p$, since $p\nmid n$; say $p\nmid x$, so $\partial F/\partial x=3x^2\not\equiv0\pmod p$ (here $p\ne3$). Let $(x,y,z)$
solve $F\equiv0\pmod{p^j}$ with $j\ge1$ and $p\nmid x$. The points congruent to it modulo $p^j$ are $(x+p^ja,y+p^jb,z+p^jc)$ with $a,b,c$ modulo $p$, and by Taylor's formula (all higher terms are
divisible by $p^{2j}$, hence by $p^{j+1}$)
\[
F(x+p^ja,y+p^jb,z+p^jc)\equiv F(x,y,z)+p^j\bigl(3x^2a+3y^2b+3z^2c\bigr)\pmod{p^{j+1}}.
\]
Dividing by $p^j$, the condition $F\equiv0\pmod{p^{j+1}}$ is one linear equation in $(a,b,c)$ over $\mathbb F_p$ with a nonzero coefficient $3x^2$, which has exactly $p^2$ solutions. So each solution modulo $p^j$ has exactly $p^2$
lifts modulo $p^{j+1}$, all with $p\nmid x$ again. By induction there are $p^2\cdot(p^2)^{k-1}=p^{2k}$ solutions modulo $p^k$, and the ratio in the definition of $\sigma_p(n)$ is $1$ for every $k$.

(b) By the proof of Proposition~\ref{prop:mod9}, there is no solution modulo $9$, hence none modulo $3^k$ for $k\ge2$.
\end{proof}

For the primes $p\equiv1\pmod3$, cubing is three-to-one on $\mathbb F_p^\times$ and $\sigma_p(n)$ is not $1$. For $p\nmid3n$ the lifting argument above still applies, so $\sigma_p(n)$ is the count modulo $p$ divided by $p^2$, and the theory of Jacobi sums shows that it is $1+O(1/p)$. Numerically, one finds a striking pattern: $p(\sigma_p(n)-1)$ is close to $6$ when $n$ is a cube modulo $p$ and close to $-3$ when it is not. A random $n$ is a cube modulo $p$ with probability $\frac13$, and $\frac13\cdot6+\frac23\cdot(-3)=0$: the corrections cancel on average, and the product over $p$ converges, though only conditionally and slowly. Computing it accurately is a small research problem in itself, which Booker and Sutherland solved.

\begin{remark}[Heath-Brown's conjecture, 1992]\label{rem:hb}
Combining the two factors, Heath-Brown conjectured that for $n\not\equiv\pm4\pmod9$ and $n$ not a cube, the number of solutions with $\max(|x|,|y|,|z|)\le B$, counted up to permutation, satisfies
\[
N_n(B)\sim\rho(n)\log B,\qquad\rho(n)=\frac16\,\sigma_\infty\prod_p\sigma_p(n)>0.
\]
In particular every such $n$ would be a sum of three cubes in infinitely many ways. This is the same shape of formula that the circle method \emph{proves} for Waring's problem with
enough variables (the asymptotic for Waring's problem with many variables, proved in Volume~5); for three cubes the method is far out of reach, and the formula is only a prediction.
\end{remark}
